% !TEX root = _yoko.tex

% ==========================================
% プロジェクト予稿（1P）の本文
% ==========================================

\section{はじめに}
近年インターネットの普及に伴い音楽をインターネット配信で楽しむ人が増加している．
各音楽配信サイトなどでは，ユーザの嗜好に合った楽曲の検索を容易にするために楽曲を音楽ジャンル別で分類されることが多くみられる．

しかし，楽曲のジャンルは時系列で変わっていくのではないだろうか？

また，音楽ジャンル解析は古くから数多くの研究が行われており，近年では主要な研究テーマの１つとして扱われている．
しかし，これらの大半は１つの楽曲を1つの音楽ジャンルに分類するものである．
このような音楽ジャンルの分類方法では，楽曲内における音楽ジャンルの変化の識別が困難であると考える．

そこで本研究では，数ある音響特徴量の中でも音色に着目し，対象楽曲の楽器構成から音楽ジャンルの変化を時系列で表示するシステムを提案することを目的とする．


\section{研究概要}
入力は，MP3形式の音響信号とする．
入力音声に対し，webサイトを介した処理を行い，楽器構成を抽出する．次に，抽出された楽器構成に対してデータベースとの照合を行い，ジャンルの一致率を算出する．
分析結果はジャンルごとにグラフとして視覚化する．


\subsection{楽曲データベース}
音楽ジャンル分析には，独自に収集したデータセットを使用する．
jazz，classic，rockの３ジャンルを対象とし，各ジャンル50曲，合計150曲の楽曲を選定した．


\subsection{楽器構成の抽出}
まず，入力音声ファイルを５秒ごとのセグメントに分割する．

次に，各セグメントに対し，LP-MusicCaps\cite{LP}を用いて楽器特徴量を解析した上での楽器構成を抽出する．

抽出されたテキストのうち，楽器名の記述部分のみを抽出し，当該セグメントの楽器構成としてリスト化する．
以上の処理をすべてのセグメントに対して行い，時間ごとの楽器構成を抽出する．


\subsection{音楽ジャンルの一致率分析}
本研究では，対象楽曲の楽器構成と一致する楽器構成の楽曲がデータベースに何曲存在するかを計算した数値をジャンルの一致率と定義する．

楽曲Xにおける
ジャンル$G$の一致率$A_X(G_i)$の計算式は以下のとおりである．

\begin{equation}
	A_X(G_i) = \frac{\text{一致した楽曲数}}{\text{ジャンル}G_i\text{の総楽曲数}}
\end{equation}

$A_X(G_i)$は最大値を1.0として，数値が大きいほど一致度が高いことを示す．

算出した一致率を，横軸を時系列，縦軸を最大値１の音楽ジャンル一致率とするヒストグラムで表示する．


\section{検証}
入力する楽曲はsnowmanの「EMPIRE」冒頭1分19秒を使用する．

この楽曲はモーツァルト作曲「交響曲第25番ト短調K.183」の第一楽章をサンプリングした楽曲であり，classicとjazzの要素が時間水位に伴い混在する楽曲となっている．

対象楽曲の冒頭5秒間の楽器構成はviolin，viola，cello，contrabass，oboe，faggot，corn，horn，vocal，drumである．
そのため，classicの一致度が高くなることが考えられた．

しかし，実装結果はclassicの一致度が0となっている．
データベースのClassic楽曲にはvocalが使用されている楽曲がないため，このような結果が出ていると考えられる．

また，冒頭5秒間のセグメントにて取得した楽器構成リストは，['vocal']のみである．
対象楽曲で試用されている楽器数が多いため，テキスト生成の際に正確に楽器を抽出できなかったものと考えられる．


\section{今後の課題}
本システムでは音色の抽出にLP-MusicCapsを使用したが，楽器数が多いと正確に使用楽器が抽出できないという問題点があった．
そのため，楽器数の少ない楽曲を対象楽曲とした際は正しい結果が得られたが，音数が多く複雑な楽曲を対象楽曲とした際に正しい実装結果が得られなかった．

また，今回は50曲×3の150曲のデータベースを使用した．
しかし，より正確かつ汎用性の高いシステムになるために，データベースの拡大が必要であると考える．


\section{おわりに}
本研究では，楽曲中の時間推移に伴う音楽ジャンルの変化を分析・可視化するシステムを提案し，検証を行った．
従来の1つの楽曲を1つの音楽ジャンルに分類する解析手法では捉えきれなかったジャンルの多様性に注目し，楽器構成の変化を時間軸に沿って解析することで，楽曲内でジャンルがどのように移り変わるかを可視化する新しいアプローチを示した．

検証の結果，楽器構成がシンプルな楽曲では有効な結果を得られたが，楽器構成が複雑な楽曲では音色抽出の精度に課題があることが明らかになった．
例えば、特定の楽器が正確に認識されないず，データベースとの一致率に影響を及ぼすケースが確認された．
また，今回使用したデータベースは150曲と規模が限られており，データの多様性においても改善の余地があると考える．

今後の課題として，音色抽出アルゴリズムの改良や，より大規模で多様性に富んだデータベースの構築が挙げられる．
これにより，複雑な楽曲への対応力を向上させるとともに，システムの汎用性と精度を高めることが可能となる．







